\begin{longtable}{|>{\centering\arraybackslash}m{3.5cm}|>{\centering\arraybackslash}m{3cm}|m{7cm}|}
\caption{Legenda \textit{Sequence Diagram}}
\label{tab:legenda-sequence} \\
\hline
\textbf{Simbol} & \textbf{Nama} & \multicolumn{1}{c|}{\textbf{Keterangan}} \\
\hline
\endfirsthead
\hline
\textbf{Simbol} & \textbf{Nama} & \multicolumn{1}{c|}{\textbf{Keterangan}} \\
\hline
\endhead

\begin{minipage}[c][1.5cm][c]{3.5cm}
  \centering
  \setlength{\unitlength}{1mm}
  \begin{picture}(10,14)
    \put(4,0){\framebox(4,14){}}
  \end{picture}
\end{minipage} &
\textit{Activation} &
Periode waktu di mana suatu objek sedang aktif menjalankan operasi atau menunggu respons. \\
\hline

\begin{minipage}[c][1.5cm][c]{3.5cm}
  \centering
  \setlength{\unitlength}{1mm}
  \begin{picture}(30,14)
    \put(0,11){\framebox(30,3){\tiny Object}}
    \multiput(14,0)(0,2){5}{\line(0,1){1}}
  \end{picture}
\end{minipage} &
\textit{Object Lifeline} &
Garis vertikal putus-putus yang merepresentasikan keberadaan suatu objek atau komponen selama rentang waktu tertentu dalam interaksi. \\
\hline

\begin{minipage}[c][1.8cm][c]{3.5cm}
  \centering
  \setlength{\unitlength}{1mm}
  \begin{picture}(20,18)
    \put(10,16){\circle{3}}
    \put(10,14.5){\line(0,-1){4}}
    \put(6,12){\line(1,0){8}}
    \put(10,10.5){\line(-2,-3){2}}
    \put(10,10.5){\line(2,-3){2}}
    \multiput(10,0)(0,2){3}{\line(0,1){1}}
  \end{picture}
\end{minipage} &
\textit{Actor Lifeline} &
Garis vertikal putus-putus yang menggambarkan keberadaan seorang aktor selama interaksi berlangsung. \\
\hline

\begin{minipage}[c][1.5cm][c]{3.5cm}
  \centering
  \setlength{\unitlength}{1mm}
  \begin{picture}(30,14)
    \put(0,0){\framebox(30,14){}}
    \put(0,10){\framebox(8,4){\tiny loop}}
  \end{picture}
\end{minipage} &
\textit{Loop Fragment} &
Fragmen kombinasi yang menunjukkan blok pesan yang diulang selama kondisi tertentu terpenuhi. \\
\hline

\begin{minipage}[c][1.5cm][c]{3.5cm}
  \centering
  \setlength{\unitlength}{1mm}
  \begin{picture}(30,14)
    \put(0,0){\framebox(30,14){}}
    \put(0,10){\framebox(6,4){\tiny opt}}
  \end{picture}
\end{minipage} &
\textit{Optional Fragment} &
Fragmen kombinasi yang menunjukkan blok pesan yang hanya dieksekusi apabila kondisi penjaga (\textit{guard condition}) terpenuhi. \\
\hline

\begin{minipage}[c][2cm][c]{3.5cm}
  \centering
  \setlength{\unitlength}{1mm}
  \begin{picture}(30,18)
    \put(0,0){\framebox(30,18){}}
    \put(0,14){\framebox(6,4){\tiny alt}}
    \multiput(0,9)(2,0){15}{\line(1,0){1}}
  \end{picture}
\end{minipage} &
\textit{Alternative Fragment} &
Fragmen kombinasi yang menunjukkan dua atau lebih blok pesan alternatif, di mana hanya satu blok yang dieksekusi berdasarkan kondisi yang berlaku. \\
\hline

\begin{minipage}[c][1cm][c]{3.5cm}
  \centering
  \setlength{\unitlength}{1mm}
  \begin{picture}(30,6)
    \put(2,3){\vector(1,0){26}}
  \end{picture}
\end{minipage} &
\textit{Message} &
Panah solid yang merepresentasikan pengiriman pesan atau pemanggilan operasi dari satu objek ke objek lain. \\
\hline

\begin{minipage}[c][1cm][c]{3.5cm}
  \centering
  \setlength{\unitlength}{1mm}
  \begin{picture}(30,6)
    \multiput(2,3)(3,0){9}{\line(1,0){2}}
    \put(28,3){\vector(1,0){0}}
  \end{picture}
\end{minipage} &
\textit{Return Message} &
Panah putus-putus yang merepresentasikan pengembalian nilai atau respons dari operasi yang telah dieksekusi. \\
\hline

\end{longtable}